\documentclass[10pt,a4paper]{amsart}
\usepackage{amsmath,amssymb,amsthm,graphicx}
% Retained minimal preamble for the Stacks Project chapter sites-cohomology.tex (Cohomology on Sites).
% Verbatim excerpt of the chapter preamble: Xy-pic setup (native preamble lines 14-18) and the native
% theorem declarations with the mathematical macros (native preamble lines 154-217). Copied mechanically;
% nothing added, removed or redefined. The wrapper class and the amsmath/amssymb/amsthm/graphicx base
% packages are supplied by the assembler, so the class/externaldocument part of the native preamble is omitted.
% For commutative diagrams we use Xy-pic
\usepackage[all]{xy}

% We use 2cell for 2-commutative diagrams.
\xyoption{2cell}
% Theorem environments.
%
\theoremstyle{plain}
\newtheorem{theorem}[subsection]{Theorem}
\newtheorem{proposition}[subsection]{Proposition}
\newtheorem{lemma}[subsection]{Lemma}

\theoremstyle{definition}
\newtheorem{definition}[subsection]{Definition}
\newtheorem{example}[subsection]{Example}
\newtheorem{exercise}[subsection]{Exercise}
\newtheorem{situation}[subsection]{Situation}

\theoremstyle{remark}
\newtheorem{remark}[subsection]{Remark}
\newtheorem{remarks}[subsection]{Remarks}

\numberwithin{equation}{subsection}

% Macros
%
\def\lim{\mathop{\mathrm{lim}}\nolimits}
\def\colim{\mathop{\mathrm{colim}}\nolimits}
\def\Spec{\mathop{\mathrm{Spec}}}
\def\Hom{\mathop{\mathrm{Hom}}\nolimits}
\def\Ext{\mathop{\mathrm{Ext}}\nolimits}
\def\SheafHom{\mathop{\mathcal{H}\!\mathit{om}}\nolimits}
\def\SheafExt{\mathop{\mathcal{E}\!\mathit{xt}}\nolimits}
\def\Sch{\mathit{Sch}}
\def\Mor{\mathop{\mathrm{Mor}}\nolimits}
\def\Ob{\mathop{\mathrm{Ob}}\nolimits}
\def\Sh{\mathop{\mathit{Sh}}\nolimits}
\def\NL{\mathop{N\!L}\nolimits}
\def\CH{\mathop{\mathrm{CH}}\nolimits}
\def\proetale{{pro\text{-}\acute{e}tale}}
\def\etale{{\acute{e}tale}}
\def\QCoh{\mathit{QCoh}}
\def\Ker{\mathop{\mathrm{Ker}}}
\def\Im{\mathop{\mathrm{Im}}}
\def\Coker{\mathop{\mathrm{Coker}}}
\def\Coim{\mathop{\mathrm{Coim}}}

% Boxtimes
%
\DeclareMathSymbol{\boxtimes}{\mathbin}{AMSa}{"02}

%
% Macros for moduli stacks/spaces
%
\def\QCohstack{\mathcal{QC}\!\mathit{oh}}
\def\Cohstack{\mathcal{C}\!\mathit{oh}}
\def\Spacesstack{\mathcal{S}\!\mathit{paces}}
\def\Quotfunctor{\mathrm{Quot}}
\def\Hilbfunctor{\mathrm{Hilb}}
\def\Curvesstack{\mathcal{C}\!\mathit{urves}}
\def\Polarizedstack{\mathcal{P}\!\mathit{olarized}}
\def\Complexesstack{\mathcal{C}\!\mathit{omplexes}}
% \Pic is the operator that assigns to X its picard group, usage \Pic(X)
% \Picardstack_{X/B} denotes the Picard stack of X over B
% \Picardfunctor_{X/B} denotes the Picard functor of X over B
\def\Pic{\mathop{\mathrm{Pic}}\nolimits}
\def\Picardstack{\mathcal{P}\!\mathit{ic}}
\def\Picardfunctor{\mathrm{Pic}}
\def\Deformationcategory{\mathcal{D}\!\mathit{ef}}

% --- appended 2026-10-12 (Chief): the project own definitions of the reference/slogan/history
% display environments, copied mechanically from the shared Stacks preamble lines 8-12
% (\usepackage{verbatim} plus the three \newenvironment lines, which the chapter preamble keeps
% verbatim from the project preamble). Nothing added, removed or redefined: these environments
% carry no mathematical content and the project itself typesets them as comments. Without these
% three lines a counted statement whose own text uses them fails with Environment slogan undefined.
\usepackage{verbatim}
\newenvironment{reference}{\comment}{\endcomment}
\newenvironment{slogan}{\comment}{\endcomment}
\newenvironment{history}{\comment}{\endcomment}
\usepackage{cleveref}
\crefname{originalreference}{original reference}{original references}
\Crefname{originalreference}{Original reference}{Original references}
\makeatletter
\newlabel{lemma-Leray}{{01F2}{source}{}{sourcepointer-1}{}}
\newlabel{lemma-Leray@cref}{{[originalreference][0][]01F2}{[1][0][]source}{}{}{}}
\makeatother
\crefname{originalreference}{original reference}{original references}
\Crefname{originalreference}{Original reference}{Original references}
\makeatletter
\newlabel{situation-compare}{{0EZ3}{source}{}{sourcepointer-2}{}}
\newlabel{situation-compare@cref}{{[originalreference][0][]0EZ3}{[1][0][]source}{}{}{}}
\makeatother
\title{Under condition V the pullback on cohomology is injective with image the classes trivialised by a covering}
\author{The Stacks Project Authors}
\date{stacks-project:sites-cohomology.tex:local-snapshot-20261011; GNU FDL 1.2 or later}
\begin{document}
\maketitle
\noindent\textit{Editorial note.} All delivered mathematical text is the author's own, copied from the pinned Stacks Project chapter file sites-cohomology.tex; no statement or proof step was rewritten, reordered or removed.\par\medskip
\noindent\textit{Editorial note.} The counted proof is the author's complete argument: the proof is written out in full: the reformulation of the map through epsilon\_X, the exact sequence obtained from the Leray spectral sequence, and the reading off of injectivity and the image.\par\medskip
\noindent\textit{Editorial note.} The label lemma-V-implies-cohomology-extra-general names the counted statement, and the proof environment immediately following it in the source is delivered verbatim as its proof.\par\medskip
\noindent\textit{Editorial note.} The counted statement is carried by the author's own lemma environment in the chapter source.\par\medskip
\noindent\textit{Editorial note.} The statement refers to the standing situation of the comparison of two topologies and is delivered verbatim; the two cited inputs are kept as external pointers with their canonical Stacks tags.\par\medskip
\noindent\textit{Allowed standard prerequisites (editorial).}\begin{itemize}
\item The standing situation comparing two topologies and the functors epsilon\_X, epsilon\_X\textasciicircum{}\{-1\}. Same chapter, the situation comparing the topologies tau and tau' together with the subcategories A\_X and the quadratic identity epsilon\_X\textasciicircum{}\{-1\} epsilon\_\{X,*\} = id on A\_X. Supplies the objects and the description of the map as pullback of cohomology classes.
\item The Leray spectral sequence for the morphism epsilon\_X. Same chapter, the Leray spectral sequence E\_2\textasciicircum{}\{p,q\} = H\textasciicircum{}p\_\{tau'\}(X, R\textasciicircum{}q epsilon\_\{X,*\}F) converging to H\textasciicircum{}\{p+q\}\_tau(X, F). Supplies the exact sequence from which the statement is read off.
\end{itemize}\medskip
\ifdefined\hypertarget\hypertarget{sourcepointer-1}{}\fi
\noindent\textit{Source reference.} Original reference 01F2: sources/stacks-sites-cohomology/pinned-original-source.tex, source lines 2323--2323. Prerequisite pointer invoked at proof line(s) 7073; it names a standard fact of the same chapter that the counted argument uses, and every step of the counted proof built on it is retained verbatim in the delivered body.\par\medskip
\ifdefined\hypertarget\hypertarget{sourcepointer-2}{}\fi
\noindent\textit{Source reference.} Original reference 0EZ3: sources/stacks-sites-cohomology/pinned-original-source.tex, source lines 6841--6841. Prerequisite pointer invoked at proof line(s) 7062; it names a standard fact of the same chapter that the counted argument uses, and every step of the counted proof built on it is retained verbatim in the delivered body.\par\medskip
\begin{lemma}
\label{lemma-V-implies-cohomology-extra-general}
In Situation \ref{situation-compare} if $(V_n)$ holds, then for
$X$ in $\mathcal{C}$ and $\mathcal{F}$ in $\mathcal{A}_X$ the map
$H^{n + 1}_{\tau'}(X, \epsilon_{X, *}\mathcal{F}) \to
H^{n + 1}_\tau(X, \mathcal{F})$
is injective with image those classes which become trivial on
a $\tau'$-covering of $X$.
\end{lemma}

\begin{proof}
Recall that $\epsilon_X^{-1}\epsilon_{X, *}\mathcal{F} = \mathcal{F}$
hence the map is given by pulling back cohomology classes
by $\epsilon_X$. The Leray spectral sequence (Lemma \ref{lemma-Leray})
$$
E_2^{p, q} = H^p_{\tau'}(X, R^q\epsilon_{X, *}\mathcal{F})
\Rightarrow
H^{p + q}_\tau(X, \mathcal{F})
$$
combined with the assumed vanishing gives an exact sequence
$$
0 \to
H^{n + 1}_{\tau'}(X, \epsilon_{X, *}\mathcal{F}) \to
H^{n + 1}_\tau(X, \mathcal{F}) \to
H^0_{\tau'}(X, R^{n + 1}\epsilon_{X, *}\mathcal{F})
$$
This is a restatement of the lemma.
\end{proof}

\end{document}
